\clearpage\thispagestyle{empty}
\pdfbookmark[0]{Quellen zu Beobachtung (7. März 1812)}{observation1812march7-sources}
\begin{center}\textbf{Quellen zu Beobachtung (7. März 1812)}\end{center}
\noindent Diese KI-Übertragung ins heutige Deutsch folgt den \href{https://gdz.sub.uni-goettingen.de/id/PPN236007467}{\textit{Werken}, Band~VI} (Göttingen, 1874), S.~348. Der Bericht erschien unter der Überschrift Göttingen in den \textit{Göttingischen gelehrten Anzeigen}, Heft~38 vom 7. März 1812, S.~\href{https://api.digitale-sammlungen.de/iiif/image/v2/bsb10538681_00413/full/1600,/0/default.jpg}{369} und \href{https://api.digitale-sammlungen.de/iiif/image/v2/bsb10538681_00414/full/1600,/0/default.jpg}{370}. Das digitalisierte Exemplar stammt aus der Bayerischen Staatsbibliothek.
\par\smallskip\noindent Beim ersten Eintrag zu Mayer steht in der Werkausgabe 16c und in der Zeitschrift auf S.~369 die Zahl 160. Für 162 Mayer gibt dieser Bericht 27.0 Sekunden für Gauß und 27.5 für Harding an. Im Brief vom 1. März ordnet Gauß die Werte umgekehrt zu. KI-Übertragung und KI-Satz.
\clearpage
